% Cayley's Collected Mathematical Papers, Volume 9, Pages 1-50
% Transcribed from: collectedmathema09cayluoft.pdf

\documentclass[12pt,a4paper]{article}

% Standard English mathematical preamble
\usepackage[utf8]{inputenc}
\usepackage[T1]{fontenc}
\usepackage[english]{babel}
\usepackage{amsmath,amssymb,amsthm}
\usepackage{geometry}
\usepackage{titlesec}
\usepackage{fancyhdr}
\usepackage{enumitem}
\usepackage{array}
\usepackage{booktabs}
\usepackage{longtable}
\usepackage{multicol}
\usepackage{hyperref}

% Geometry settings
\geometry{top=1in,bottom=1in,left=1.25in,right=1.25in}

% Header/Footer
\pagestyle{fancy}
\fancyhf{}
\rhead{Cayley, Collected Math. Papers, Vol. IX}
\lhead{Pages 1--50}
\rfoot{Page \thepage}

% Theorem-like environments
\theoremstyle{plain}
\newtheorem{theorem}{Theorem}
\newtheorem{lemma}[theorem]{Lemma}
\newtheorem{corollary}[theorem]{Corollary}
\theoremstyle{definition}
\newtheorem{definition}{Definition}

% Document metadata
\title{\textbf{THE COLLECTED MATHEMATICAL PAPERS}\\[6pt] \large OF\\[6pt] \textbf{ARTHUR CAYLEY, Sc.D., F.R.S.,}\\[4pt] \footnotesize LATE SADLERIAN PROFESSOR OF PURE MATHEMATICS IN THE UNIVERSITY OF CAMBRIDGE.\\[12pt] \large \textbf{VOL. IX.}}
\author{}
\date{}

\begin{document}

% =====================================================================
% TITLE PAGE (Pages 1-8 of PDF)
% =====================================================================

\thispagestyle{empty}

\begin{center}

\vspace*{2cm}

{\Large FOR}\\[8pt]
{\large m}\\[8pt]
{\Large IN}\\[8pt]
{\Large LIBRARY}\\[8pt]
{\Large OHIY}\\[12pt]
{\large SEEN BY}\\[8pt]
{\Large PRESERVATION}\\[8pt]
{\Large SERVICES}\\[8pt]
{\large DATE}\\[16pt]
{\Large C'\^{}-\^{},\^{}}

\vfill
\newpage

{\Large MATHEMATICAL PAPERS.}

\vfill
\newpage

{\Large london: C. J. CLAY and SONS,}\\[6pt]
\normalsize CAMBRIDGE UNIVERSITY PRESS WAREHOUSE,\\
AVE MARIA LANE.\\[6pt]
Glasgow: 263, ARGYLE STREET.\\[6pt]
Leipzig: F. A. BROCKHAUS.\\[6pt]
New York: MACMILLAN AND CO.

\vfill
\newpage

\vspace*{3cm}

{\Large \textbf{THE COLLECTED}}\\[8pt]
{\Large \textbf{MATHEMATICAL PAPERS}}\\[12pt]
{\large OF}\\[12pt]
{\Large \textbf{ARTHUR CAYLEY, Sc.D., F.R.S.,}}\\[8pt]
\normalsize LATE SADLERIAN PROFESSOR OF PURE MATHEMATICS\\
IN THE UNIVERSITY OF CAMBRIDGE.\\[20pt]
{\Large \textbf{VOL. IX.}}\\[30pt]
{\Large CAMBRIDGE:}\\[6pt]
{\large AT THE UNIVERSITY PRESS.}\\[12pt]
{\small All Rights reserved.}

\vfill
\newpage

\vspace*{6cm}

{\large CAMBRIDGE:}\\[6pt]
PRINTED BY J. AND C. F. CLAY,\\
AT THE UNIVERSITY PRESS.

\end{center}

% =====================================================================
% ADVERTISEMENT (Page 9 of PDF)
% =====================================================================

\newpage
\thispagestyle{empty}

\begin{center}
{\large \textbf{ADVERTISEMENT.}}
\end{center}

\vspace{12pt}

\noindent THE present volume contains 74 papers, numbered 556 to 629, published for the most part in the years 1874 to 1877.

\vspace{12pt}

\noindent The Table for the nine volumes is

\begin{center}
\begin{tabular}{r @{~} l @{~} r @{~} l}
Vol. & I. & Numbers & 1 to 100,\\ 
`` & II. & `` & 101 ,, 158,\\ 
`` & III. & `` & 159 ,, 222,\\ 
`` & IV. & `` & 223 ,, 299,\\ 
`` & V. & `` & 300 ,, 383,\\ 
`` & VI. & `` & 384 ,, 416,\\ 
`` & VII. & `` & 417 ,, 485,\\ 
`` & VIII. & `` & 486 ,, 555,\\ 
`` & IX. & `` & 556 ,, 629.\\ 
\end{tabular}
\end{center}

\vspace{12pt}

\hfill A. E. Forsyth.

\hfill 17 December 1895.

% =====================================================================
% DIGITIZATION NOTICE (Page 10 of PDF)
% =====================================================================

\newpage
\thispagestyle{empty}

\begin{center}
{\small Digitized by the Internet Archive\\
in 2007 with funding from\\
Microsoft Corporation}

\vspace{12pt}
\url{http://www.archive.org/details/collectedmathema09cayluoft}
\end{center}

% =====================================================================
% TABLE OF CONTENTS (Pages 11-17 of PDF)
% =====================================================================

\newpage
\setcounter{page}{7}

\begin{center}
{\large \textbf{CONTENTS.}}
\end{center}

\vspace{12pt}

\noindent
\begin{longtable}{r @{. } p{4.8in} r}
556 & On Steiner's surface & 1 \\
  & Proc.\ Lond.\ Math.\ Society, t.\ v.\ (1873--1874), pp.\ 14--25 & \\ 

557 & On certain constructions for bicircular quartics & 13 \\
  & Proc.\ Lond.\ Math.\ Society, t.\ v.\ (1873--1874), pp.\ 29--31 & \\ 

558 & A geometrical interpretation of the equations obtained by equating to zero the resultant and the discriminants of two binary quantics & 16 \\
  & Proc.\ Lond.\ Math.\ Society, t.\ v.\ (1873--1874), pp.\ 31--33 & \\ 

559 & [Note on inversion] & 18 \\
  & Proc.\ Lond.\ Math.\ Society, t.\ v.\ (1873--1874), p.\ 112 & \\ 

560 & [Addition to Lord Rayleigh's paper ``On the numerical calculation of the roots of fluctuating functions''] & 19 \\
  & Proc.\ Lond.\ Math.\ Society, t.\ v.\ (1873--1874), pp.\ 123, 124 & \\ 

561 & On the geometrical representation of Cauchy's theorems of root-limitation & 21 \\
  & Camb.\ Phil.\ Trans., t.\ xii.\ Part ii.\ (1877), pp.\ 395--413 & \\ 

562 & On a theorem in maxima and minima & 40 \\
  & Quart.\ Math.\ Jour., t.\ x.\ (1870), pp.\ 262, 263 & \\ 

563 & On a theorem in elimination & 42 \\
  & Quart.\ Math.\ Jour., t.\ xii.\ (1873), pp.\ 5, 6 & \\ 

564 & Note on the transformation of two simultaneous equations & 43 \\
  & Quart.\ Math.\ Jour., t.\ xii.\ (1873), pp.\ 16--19 & \\ 

565 & Note on the Cartesian & 45 \\
  & Quart.\ Math.\ Jour., t.\ xii.\ (1873), pp.\ 16--19 & \\ 

566 & On the transformation of the equation of a surface to a set of chief axes & 48 \\
  & Quart.\ Math.\ Jour., t.\ xii.\ (1873), pp.\ 34--38 & \\ 

567 & On an identical equation connected with the theory of invariants & 52 \\
  & Quart.\ Math.\ Jour., t.\ xii.\ (1873), pp.\ 115--118 & \\ 

568 & Note on the integrals $\displaystyle\int_0^\infty \cos x^n \,dx$ and $\displaystyle\int_0^\infty \sin x^n \,dx$ & 56 \\
  & Quart.\ Math.\ Jour., t.\ xii.\ (1873), pp.\ 118--126 & \\ 

569 & On the cyclide & 64 \\
  & Quart.\ Math.\ Jour., t.\ xii.\ (1873), pp.\ 148--165 & \\ 

570 & On the superlines of a quadric surface in five-dimensional space & 79 \\
  & Quart.\ Math.\ Jour., t.\ xii.\ (1873), pp.\ 185--191 & \\ 

571 & A demonstration of Dupin's theorem & 84 \\
  & Quart.\ Math.\ Jour., t.\ xii.\ (1873), pp.\ 193--197 & \\ 

572 & On the condition for the existence of a given set of lines are the same quartic functions of $x$, $y$ respectively & 90 \\
  & Quart.\ Math.\ Jour., t.\ xii.\ (1873), pp.\ 197, 198 & \\ 

573 & On an algebraical operation & 94 \\
  & Quart.\ Math.\ Jour., t.\ xii.\ (1873), pp.\ 221--228 & \\ 

574 & On Wronski's theorem & 96 \\
  & Quart.\ Math.\ Jour., t.\ xii.\ (1873), pp.\ 176--180 & \\ 

575 & On a special quartic transformation of an elliptic function & 103 \\
  & Quart.\ Math.\ Jour., t.\ xii.\ (1873), pp.\ 266--269 & \\ 

576 & Addition to Mr Walton's paper ``On the ray-planes in biaxal crystals'' & 107 \\
  & Quart.\ Math.\ Jour., t.\ xii.\ (1873), pp.\ 273--275 & \\ 

577 & Note in illustration of certain general theorems obtained by Dr Lipschitz & 110 \\
  & Quart.\ Math.\ Jour., t.\ xii.\ (1873), pp.\ 346--349 & \\ 
\end{longtable}

% =====================================================================
% CONTINUED TABLE OF CONTENTS (Pages 12-17)
% =====================================================================

\newpage

\begin{center}
{\large \textbf{CONTENTS.}}
\end{center}

\vspace{12pt}

\noindent
\begin{longtable}{r @{. } p{4.8in} r}
578 & A memoir on the transformation of elliptic functions & 116 \\
  & Phil.\ Trans., t.\ CLXIV.\ (for 1874), pp.\ 397--456 & \\ 

579 & Address delivered by the President, Professor Cayley, on presenting the Gold Medal of the [Royal Astronomical] Society to Professor Simon Newcomb & 176 \\
  & Monthly Notices R.\ Ast.\ Society, t.\ xxxiv.\ (1873--1874), pp.\ 224--233 & \\ 

580 & On the number of distinct terms in a symmetrical or partially symmetrical determinant; with an addition & 185 \\
  & Monthly Notices R.\ Ast.\ Society, t.\ xxxiv.\ (1873--1874), pp.\ 303--307; p.\ 335 & \\ 

581 & On a theorem in elliptic motion & 191 \\
  & Monthly Notices R.\ Ast.\ Society, t.\ xxxv.\ (1874--1875), pp.\ 337--339 & \\ 

582 & Note on the Theory of Precession and Nutation & 194 \\
  & Monthly Notices R.\ Ast.\ Society, t.\ xxxv.\ (1874--1875), pp.\ 340--343 & \\ 

583 & On spheroidal trigonometry & 197 \\
  & Monthly Notices R.\ Ast.\ Society, t.\ xxxvii.\ (1876--1877), p.\ 92 & \\ 

584 & Addition to Prof.\ R.\ S.\ Ball's paper ``Note on a transformation of Lagrange's equations of motion in generalised coordinates, which is convenient in Physical Astronomy'' & 198 \\
  & Monthly Notices R.\ Ast.\ Society, t.\ xxxvii.\ (1876--1877), pp.\ 269--271 & \\ 

585 & A new theorem on the equilibrium of four forces acting on a solid body & 201 \\
  & Phil.\ Mag., t.\ XXXI.\ (1866), pp.\ 78, 79; Camb.\ Phil.\ Soc.\ Proc., t.\ I.\ (1866), p.\ 235 & \\ 

586 & On the mathematical theory of isomers & 202 \\
  & Phil.\ Mag., t.\ XLVII.\ (1874), pp.\ 444--467 & \\ 

587 & A Smith's Prize dissertation [1873] & 205 \\
  & Messenger of Mathematics, t.\ iii.\ (1874), pp.\ 1-- & \\ 

588 & Problem, [on tetrahedra] & 209 \\
  & Messenger of Mathematics, t.\ iii.\ (1874), pp.\ 50--52 & \\ 

589 & On residuation in regard to a cubic curve & 211 \\
  & Messenger of Mathematics, t.\ iii.\ (1874), pp.\ 62--65 & \\ 

590 & Addition to Prof.\ Hall's paper ``On the motion of a particle toward an attracting centre at which the force is infinite.'' & 215 \\
  & Messenger of Mathematics, t.\ iii.\ (1874), pp.\ 149--152 & \\ 

591 & A Smith's Prize paper and dissertation [1874]; solutions and remarks & 218 \\
  & Messenger of Mathematics, t.\ iii.\ (1874), pp.\ 165--183; t.\ iv.\ (1875), pp.\ 6--8 & \\ 

592 & On the Mercator's projection of a skew hyperboloid of revolution & 237 \\
  & Messenger of Mathematics, t.\ iii.\ (1874), pp.\ 187, 188 & \\ 

593 & A Smith's Prize paper and dissertation [1874]; solutions and remarks & 241 \\
  & Messenger of Mathematics, t.\ iv.\ (1875), pp.\ 17-- & \\ 

594 & On a differential equation in the theory of elliptic functions & 244 \\
  & Messenger of Mathematics, t.\ iv.\ (1875), pp.\ 34--36 & \\ 

595 & On a differential equation in the theory of elliptic functions & 246 \\
  & Messenger of Mathematics, t.\ iv.\ (1875), pp.\ 69, 70 & \\ 

596 & On a Senate-House problem & 250 \\
  & Messenger of Mathematics, t.\ iv.\ (1875), pp.\ 75--78 & \\ 

597 & Note on a theorem of Jacobi's for the transformation of a double integral & 253 \\
  & Messenger of Mathematics, t.\ iv.\ (1875), pp.\ 92--94 & \\ 

598 & Note on a process of integration & 257 \\
  & Messenger of Mathematics, t.\ iv.\ (1875), pp.\ 110--113 & \\ 

599 & A Sheepshanks' problem (1866) & 259 \\
  & Messenger of Mathematics, t.\ iv.\ (1875), pp.\ 149, 150 & \\ 

600 & Theorem on the $n$-th roots of unity & 263 \\
  & Messenger of Mathematics, t.\ iv.\ (1875), pp.\ 157--160 & \\ 

601 & Note on the Cassinian & 264 \\
  & Messenger of Mathematics, t.\ iv.\ (1875), pp.\ 160-- & \\ 

602 & On the potentials of polygons and polyhedra & 266 \\
  & Proc.\ Lond.\ Math.\ Society, t.\ vi.\ (1874--1875), pp.\ 20--34 & \\ 

603 & On the potential of the ellipse and the circle & 281 \\
  & Proc.\ Lond.\ Math.\ Society, t.\ vi.\ (1874--1875), pp.\ 38--58 & \\ 

604 & Determination of the attraction of an ellipsoidal shell on an exterior point & 302 \\
  & Proc.\ Lond.\ Math.\ Society, t.\ vi.\ (1874--1875), pp.\ 58--67 & \\ 

605 & Note on a point in the theory of attraction & 312 \\
  & Proc.\ Lond.\ Math.\ Society, t.\ vi.\ (1874--1875), pp.\ 79--81 & \\ 

606 & On an expression of the coordinates of a point of a quartic curve as functions of a parameter & 315 \\
  & Proc.\ Lond.\ Math.\ Society, t.\ vi.\ (1874--1875), pp.\ 81--83 & \\ 

607 & Note on a point in the theory of attraction & 318 \\
  & Brit.\ Assoc.\ Report, 1873, pp.\ 3, 4 & \\ 

608 & [Extract from a] Report on Mathematical Tables & 424 \\
  & Brit.\ Assoc.\ Report, 1875, pp.\ 281--305 & \\ 

609 & On the analytical forms called factions & 427 \\
  & Brit.\ Assoc.\ Report, 1875, pp.\ 305--336 & \\ 

610 & Trees, with application to the theory of chemical combinations & 461 \\
  & Brit.\ Assoc.\ Report, 1875, pp.\ 257--305 & \\ 

611 & On the potential of the ellipse and the circle & 500 \\
  & Brit.\ Assoc.\ Report, 1875, Notices of Communications to the Sections, pp.\ 1624--1629 & \\ 

612 & Note sur une formule d'integration indefinie & 504 \\
  & Comptes Rendus, t.\ LXXVIII.\ (1874), pp.\ 500 & \\ 

613 & On the group of points $G_6$ on a sextic curve with five double points & 508 \\
  & Math.\ Ann., t.\ viii.\ (1875), pp.\ 359--362 & \\ 

614 & On a problem of projection & 519 \\
  & Quart.\ Math.\ Jour., t.\ xiii.\ (1875), pp.\ 19--29 & \\ 

615 & On the conic torus & 522 \\
  & Quart.\ Math.\ Jour., t.\ xiii.\ (1875), pp.\ 127--129 & \\ 

616 & A geometrical illustration of the cubic transformation in elliptic functions & 527 \\
  & Quart.\ Math.\ Jour., t.\ xiii.\ (1875), pp.\ 211--216 & \\ 

617 & On the scalene transformation of a plane curve & 535 \\
  & Quart.\ Math.\ Jour., t.\ xiii.\ (1875), pp.\ 321--328 & \\ 

618 & On the mechanical description of a Cartesian & 537 \\
  & Quart.\ Math.\ Jour., t.\ xiii.\ (1875), pp.\ 328--330 & \\ 

619 & On an algebraical operation & 543 \\
  & Quart.\ Math.\ Jour., t.\ xiii.\ (1875), pp.\ 35--38 & \\ 

620 & On the condition for the existence of a given set of lines are the same quartic functions of $x$, $y$ respectively & 546 \\
  & Quart.\ Math.\ Jour., t.\ xiii.\ (1875), pp.\ 211--216 & \\ 

621 & Correction of two numerical errors in Sohnke's paper respecting modular equations & 551 \\
  & Crelle, t.\ LXXXI.\ (1876), pp.\ 369--375 & \\ 

622 & On the number of the univalent radicals $C_n H_{2n+1}$ & 581 \\
  & Phil.\ Mag., Ser.\ 5, t.\ vii.\ (1876), pp.\ 136--166 & \\ 

623 & On three-bar motion & 587 \\
  & Proc.\ Lond.\ Math.\ Society, t.\ vii.\ (1876), pp.\ 136--166 & \\ 

624 & On a system of equations connected with Malfatti's problem & 592 \\
  & Proc.\ Lond.\ Math.\ Society, t.\ vii.\ (1876), pp.\ 38--42 & \\ 

625 & On the bicursal sextic & 609 \\
  & Proc.\ Lond.\ Math.\ Society, t.\ vii.\ (1876), pp.\ 321--328 & \\ 

626 & On the general differential equation $\frac{\partial^2 Y}{\partial x^2} + \frac{\partial^2 Y}{\partial y^2} = 0$, where $X$, $Y$ are the same quartic functions of $x$, $y$ respectively & 612 \\
  & Proc.\ Lond.\ Math.\ Society, t.\ viii.\ (1877), pp.\ 34, 35 & \\ 

627 & Geometrical illustration of a theorem relating to an irrational function of an imaginary variable & 618 \\
  & Proc.\ Lond.\ Math.\ Society, t.\ viii.\ (1877), pp.\ 212--214 & \\ 

628 & On the circular relation of Moebius & 618 \\
  & Proc.\ Lond.\ Math.\ Society, t.\ viii.\ (1877), pp.\ 220--226 & \\ 

629 & On the linear transformation of the integral $\int dt / \sqrt{T}$ & 618 \\
  & Proc.\ Lond.\ Math.\ Society, t.\ viii.\ (1877), pp.\ 226--229 & \\ 
  & Plate to face p.\ 460 & \\ 
\end{longtable}

% =====================================================================
% CLASSIFICATION (Pages 19-20 of PDF)
% =====================================================================

\newpage
\setcounter{page}{15}

\begin{center}
{\large \textbf{CLASSIFICATION.}}
\end{center}

\vspace{12pt}

\noindent \textbf{Geometry:}

\begin{multicols}{2}
Moebius' circular relation, 628\\
Malfatti's problem, 622\\
Inversion, 559\\
Projections, 592, 614\\
Mechanical construction of curves, 557, 618\\
Quartic curve and functions of a single parameter, 606\\
Cartesian, 565\\
Cassinian, 601\\
Sextic curves, 613, 624\\
Correspondence, residuation and transformation, 566, 573, 589, 617\\
Conformal representation, 627\\
Illustrations of algebraical theorems, 558, 561\\
Three-bar motion, 623\\
Spheroidal Trigonometry, 583\\
Dupin's theorem, 571\\
Cyclide, 569\\
Steiner's surface, 556\\
Surface orthogonal to set of lines, 625\\
Conic torus, 615\\
Quadrics in hyperdimensional space, 570\\
\end{multicols}

\vspace{12pt}

\noindent \textbf{Astronomy and Dynamics:}

\begin{multicols}{2}
Presidential address, 579\\
Elliptic motion, 581\\
Precession and nutation, 582\\
General equations of dynamics, 577, 584, 590\\
Potentials, attractions and prepotentials, 602, 603, 604, 605, 607\\
\end{multicols}

\newpage
\setcounter{page}{16}

\begin{center}
{\large \textbf{CLASSIFICATION.}}
\end{center}

\vspace{12pt}

\noindent \textbf{Analysis:}

\begin{multicols}{2}
Determinant, number of terms in, 580\\
Elimination, 564\\
Invariants and covariants, 567, 572, 619\\
Mathematical tables, 608, 611\\
Trees, and their applications to chemistry, 586, 610, 621\\
Roots of Unity, 600\\
Transformation of equations, 563\\
Maxima and minima, 562\\
Factions, 609\\
Fluctuating functions, 560\\
Wronski's theorem of expansion, 574\\
Integration and definite integrals, 568, 596, 598, 612\\
Elliptic functions, 626, 629\\
Transformation of elliptic functions, 575, 578, 594, 597, 616, 620\\
Smith's Prize Dissertations and Solutions, 587, 591, 599\\
Miscellaneous, 576, 585, 588, 593, 595\\
\end{multicols}

% =====================================================================
% PAPER 556: ON STEINER'S SURFACE (Pages 21-32 of PDF)
% =====================================================================

\newpage
\setcounter{page}{1}

\begin{center}
[556\\
]
\end{center}

\section*{\textbf{556.}}
\addcontentsline{toc}{section}{556. On Steiner's Surface}

\begin{center}
{\large \textbf{ON STEINER'S SURFACE.}}
\end{center}

\begin{flushright}
{\small [From the Proceedings of the London Mathematical Society, vol.\ v.\ (1873--1874), pp.\ 14--25.\ Read December 11, 1873.]}
\end{flushright}

I have constructed a model and drawings of the symmetrical form of Steiner's Surface, viz.\ that wherein the four singular tangent planes form a regular tetrahedron, and consequently the three nodal lines (being the lines joining the mid-points of opposite edges) a system of rectangular axes at the centre of the tetrahedron.

Before going into the analytical theory, I describe as follows the general form of the surface.

Take the tetrahedron, and inscribe in each face a circle (there will be, of course, two circles on each face touching at the mid-point of each edge, and having their centres on the lines joining the mid-point to the opposite vertex). The three circles on any face contain, the central circle within each and the three marginal circles outside each. Now truncate the tetrahedron, reducing the altitudes to three-fourths their original values; each circle will then be on a new face. Produce the three lines joining the mid-points of the edges beyond the centre, so as each to meet two of the new faces; and from the centre of each new face excavate downwards by planes passing through these lines, each circle will be scooped out so as to become the bounding curve of an excavated cavity. The surface will thus consist of four excavations bounded each of them by a trinodal quartic, and united together so as to form a continuous surface.

We may imagine the tetrahedron placed in two different positions, (1) resting with one of its faces on the horizontal plane, (2) with two opposite edges horizontal, or say with the horizontal plane passing through the centre of the tetrahedron and being parallel to two opposite edges.

\textbf{1.}\ In the first position, the horizontal plane meets the surface in a trinodal quartic.

This is at once seen to be the inverse of a circular cubic with three real asymptotes. In fact, the general equation of a circular cubic is $y^2 + (x - \alpha)(x - \beta)(x - \gamma) = 0$; and if the origin is on the curve, then taking $\rho^2 = x^2 + y^2$ and inverting in regard to the origin, we obtain the trinodal quartic $\rho^2 y^2 + \rho^2(x - \alpha)(x - \beta)(x - \gamma) = 0$, viz.\ $(x^2 + y^2)^2 + (x - \alpha)(x - \beta)(x - \gamma) = 0$, or $x^2 + y^2 + \sqrt{(x - \alpha)(x - \beta)(x - \gamma)} = 0$, a trinodal quartic.

But the section is also obtained directly. The equation of the surface is $\sqrt{x} + \sqrt{y} + \sqrt{z} + \sqrt{w} = 0$, where $x + y + z + w = 0$ if $x = 0$, $y = 0$, $z = 0$, $w = 0$ are the equations of the four faces of the tetrahedron. The section by the plane $w = \lambda$ is given by these two equations; and writing $X$, $Y$, $Z$ for the coordinates referred to the triangle which is the section of the tetrahedron, the equations of the section become $\sqrt{X} + \sqrt{Y} + \sqrt{Z} + \sqrt{\lambda(X + Y + Z)} = 0$.

To simplify, I write $q = \dfrac{\lambda}{2\lambda - 1}$, that is $\lambda = \dfrac{q}{2q + 1}$, the equation then is
\[
\sqrt{X} + \sqrt{Y} + \sqrt{Z} + \sqrt{(2q+1)(X + Y + Z)} = 0;
\]
or, proceeding to rationalize, we have first
\[
q(X + Y + Z) = \sqrt{YZ} + \sqrt{ZX} + \sqrt{XY},
\]
and thence
\[
q^2(X + Y + Z)^2 - (YZ + ZX + XY) = 2\sqrt{XYZ}(\sqrt{X} + \sqrt{Y} + \sqrt{Z})
\]
and finally
\[
\bigl[q^2(X + Y + Z)^2 - YZ - ZX - XY\bigr]^2 = 4(2q + 1)XYZ(X + Y + Z).
\]

This is a quartic curve, having for double tangents the four lines $X = 0$, $Y = 0$, $Z = 0$, $X + Y + Z = 0$, the last of these being the line infinity touching the curve in two imaginary points, since obviously the whole real curve lies within the triangle.

This is as it should be: the double tangents are the intersections of the plane $w = \lambda$ by the singular planes of the surface. To find the points of contact, writing for instance $Z = 0$, the equation becomes $q^2(X + Y)^2 - XY = 0$, that is
\[
\frac{Y}{X} = \frac{-1 + 2q \pm \sqrt{4q^2 - 3}}{2},
\]
giving the two points of contact equi-distant from the centre; these are imaginary if $q > \tfrac{1}{2}$, but otherwise real, which agrees with what follows. (See the Table afterwards referred to.)

The nodal lines of the surface are $(x - y = 0,\ z - w = 0)$, $(y - z = 0,\ x - w = 0)$, $(z - x = 0,\ y - w = 0)$.

\newpage
\setcounter{page}{2}

We may imagine the tetrahedron placed in two different positions, (1) resting with one of its faces on the horizontal plane, (2) with two opposite edges horizontal, or say with the horizontal plane passing through the centre of the tetrahedron and being parallel to two opposite edges; or, what is the same thing, the nodal lines form a system of rectangular axes, one of them, say that of $z$, being vertical.

And I proceed to consider, in the two cases respectively, the form of the curve of intersection by a plane parallel to the horizontal plane.

For the first position: taking $X$, $Y$, $Z$ as the coordinates referred to the triangle which is the section of the tetrahedron, we have the quartic equation in the form
\[
\bigl[q^2(X + Y + Z)^2 - YZ - ZX - XY\bigr]^2 = 4(2q + 1)XYZ(X + Y + Z).
\]

For the point in question we have $X : Y : Z = -q : -q : 2q + 1$; and taking the perpendicular from the vertex on a side as unity, the values $-q$, $-q$, $2q + 1$ will be absolute magnitudes. We thus see that the curve must have the three nodes $(2q+1, -q, -q)$, $(-q, 2q + 1, -q)$, $(-q, -q, 2q + 1)$; and it is easy to verify that this is so.

The curve will pass through the centre $X = Y = Z$, if $(\sqrt{3} - 3)^2 = 12(2q + 1)$, that is, if $3(3q - 1)^2 - 4(2q + 1) = 0$, or if $(3q + 1)^2(q - 1) = 0$.

If $q = 1$, that is, $\lambda = \tfrac{3}{4}\Lambda$, the equation is
\[
(X^2 + Y^2 + Z^2 + YZ + ZX + XY)^2 - 16XYZ(X + Y + Z) = 0,
\]
where the curve is, in fact, a pair of imaginary conics meeting in the four real points $(3, -1, -1)$, $(-1, 3, -1)$, $(-1, -1, 3)$, $(1, 1, 1)$.

To verify this, observe that, writing
\[
A = (Y - Z)(2X + Y + Z),\quad B = (Z - X)(X + 2Y + Z),\quad C = (X - Y)(X + Y + 2Z),
\]
and therefore $A + B + C = 0$, the function in $(X, Y, Z)$ is $= \tfrac{1}{2}(A^2 + B^2 + C^2)$, and thus the equation may be written in the equivalent forms each of which shows that the curve breaks up into two imaginary conics.

The foregoing value $q = 1$, or $\lambda = \tfrac{3}{4}\Lambda$, belongs to the summit or highest real point of the surface.

\newpage
\setcounter{page}{3}

In the case $3q + 1 = 0$, that is, $\lambda = \tfrac{3}{2}\lambda$, the equation is
\[
[(X + Y + Z)^2 - 4(YZ + ZX + XY)]^2 = 108XYZ(X + Y + Z),
\]
which is, in fact, the equation of a curve having the centre, or point $X = Y = Z$, for a triple point.

To verify this, write $X = \beta - \gamma + u$, $Y = \gamma - \alpha + u$, $Z = \alpha - \beta + u$; also
\[
2\Delta = (\beta - \gamma)^2 + (\gamma - \alpha)^2 + (\alpha - \beta)^2,\quad n = (\beta - \gamma)(\gamma - \alpha)(\alpha - \beta).
\]

Then we have $X + Y + Z = 3u$, $YZ + ZX + XY = 3u^2 - \Delta$, $XYZ = u^3 - u\Delta + n$, and the equation is
\begin{align*}
[9u^2 - 9(3u^2 - \Delta)]^2 - 324u(u^3 - u\Delta + n) &= 0,\\
(-2u^2 + \Delta)^2 - 4u(u^3 - u\Delta + n) &= 0,\\
\Delta^2 - 4un &= 0,
\end{align*}
where the lowest terms in $\beta - \gamma$, $\gamma - \alpha$, $\alpha - \beta$ are of the order 3, and the theorem is thus proved.

The case in question, $q = -\tfrac{1}{3}$ or $\lambda = \tfrac{1}{4}\Lambda$, is where the plane passes through the centre of the tetrahedron.

When $q = \tfrac{1}{2}$, or $\lambda = \tfrac{2}{3}\Lambda$, the equation is
\[
(X^2 + Y^2 + Z^2 - 2YZ - 2ZX - 2XY)^2 = 128XYZ(X + Y + Z).
\]

Here each of the lines $X = 0$, $Y = 0$, $Z = 0$ is an osculating tangent having with the curve a 4-pointic intersection.

When $q = 0$, or $\lambda = \tfrac{1}{2}\Lambda$, the equation is
\[
(YZ + ZX + XY)^2 - 4XYZ(X + Y + Z) = 0,
\]
that is, $Y^2Z^2 + Z^2X^2 + X^2Y^2 - 2XYZ(X + Y + Z) = 0$; viz.\ each angle of the triangle is here a cusp.

When $q = -\tfrac{1}{2}$, or $\lambda = \tfrac{1}{3}\Lambda$, the curve is
\[
[X^2 + Y^2 + Z^2 - 2(YZ + ZX + XY)]^2 = 0,
\]
viz.\ the plane is here the base of the tetrahedron, and the section is the inscribed circle taken twice.

For tracing the curves, it is convenient to find the intersections with the lines $Y - Z = 0$, $Z - X = 0$, $X - Y = 0$ drawn from the centre of the triangle to the vertices; each of these lines passes through a node, and therefore besides meets the curve in two points.

Writing, for instance, $Y = X$, the equation becomes
\[
[q^2(2X + Z)^2 - 2XZ - X^2]^2 - 4(2q + 1)X^2Z(2X + Z) = 0;
\]
viz.\ this is
\[
[qZ + (2q + 1)X]^2[q^2Z^2 + (2q^2 - 2q - 1)XZ + (4q^2 - 4q + 1)X^2] = 0,
\]
where the first factor gives the node. Equating to zero the second factor, we have
\[
[2q^2 + (2q - 1 - q)Z]^2 = Z^2[(2q - 1 - q)^2 - 4q^2 + 4(2q - 1)] = Z^2(1 - q)(1 + 2q);
\]
or, finally,
\[
q^2Z = [-2q + 1 \pm 1 \pm q\sqrt{(1 - q)(1 + 2q)}]X,
\]
giving two real values for all values of $q$ from $q = 1$ to $q = -\tfrac{1}{2}$. (See the Table afterwards referred to.)

\newpage
\setcounter{page}{4}

We may recapitulate as follows:

\smallskip

$q > 1$, or $\lambda > \tfrac{3}{4}\Lambda$; the curve is imaginary, but with three real acnodes, answering to the acnodal parts of the nodal lines:

\smallskip

$q = 1$, or $\lambda = \tfrac{3}{4}\Lambda$; the summit appears as a fourth acnode:

\smallskip

$q < 1 > \tfrac{1}{2}$, or $\tfrac{3}{4}\Lambda > \lambda > \tfrac{2}{3}\Lambda$; the curve consists of three acnodes and a trigonoid lying within the triangle and having the sides of the triangle for bitangents of imaginary contact:

\smallskip

$q = \tfrac{1}{2}$, or $\lambda = \tfrac{2}{3}\Lambda$; the curve consists of three acnodes and a trigonoid having the sides of the triangle for osculating tangents:

\smallskip

$q < \tfrac{1}{2} > 0$, or $\tfrac{2}{3}\Lambda > \lambda > \tfrac{1}{2}\Lambda$; the curve consists of three conjugate points and an indented trigonoid having the sides of the triangle for bitangents of real contact:

\smallskip

$q = 0$, or $\lambda = \tfrac{1}{2}\Lambda$; curve has the summits of the triangle for cusps:

\smallskip

$q < 0 > -\tfrac{1}{2}$, or $\tfrac{1}{2}\Lambda > \lambda > \tfrac{1}{3}\Lambda$; curve has three crunodes, or say it is a cis-centric trifolium:

\smallskip

$q = -\tfrac{1}{2}$, or $\lambda = \tfrac{1}{3}\Lambda$; curve has a triple point, or say it is a centric trifolium:

\smallskip

$q < -\tfrac{1}{2} > -\tfrac{1}{3}$, or $\tfrac{1}{3}\Lambda > \lambda > \tfrac{1}{4}\Lambda$; curve has three crunodes, or say it is a trans-centric trifolium:

\smallskip

$q = -\tfrac{1}{3}$, or $\lambda = \tfrac{1}{4}\Lambda$; curve is a two-fold circle:

\smallskip

$q < -\tfrac{1}{3}$, or $\lambda < \tfrac{1}{4}\Lambda$; the curve becomes again imaginary, consisting of three acnodes answering to the acnodal parts of the nodal lines.

\vspace{12pt}

\begin{center}
\textbf{Table of Values}
\end{center}

The Table shows the critical values. The columns are $q$, $q^2$, $q^3$, $\lambda$, and the apsidal distances and contacts for the radius vector $X = Y$, the value of $Z : X$ being calculated from the foregoing formula.

\begin{center}
\begin{tabular}{c c c c c c}
\toprule
$q$ & $q^2$ & $q^3$ & $\lambda$ & Contact, $Z=0$, $X:Y$ & Apses, $Y=X$, $X:Z=$ \\
\midrule
.75 & 1.00 & imposs. & 1 & imposs. & .032 or 7.968 \\
.70 & .666 & .6666 & .5 & imposs. & .006 or 22.44 \\
.65 & .4285 & .320 & or 3.124 & .006 & 22.44 \\
.6 & .25 & .0721 & 13.9279 & .059 & 67.941 \\
.55 & .1111 & .011 & 78.988 & .15 & 337.85 \\
.5 & 0 & or $\infty$ & .25 & $\infty$ \\
.45 & .0909 & .005 & 118.99 & .39 & 457.61 \\
.4 & .1666 & .029 & 33971 & .496 & 127.504 \\
.35 & .2308 & .060 & 1672 & .648 & 61.79 \\
.3 & .2857 & .099 & 10.151 & .812 & 37.187 \\
.25 & .3333 & .141 & 6.854 & 1.25 & 25 \\
.2 & .375 & .207 & 4.904 & 1.218 & 17.89 \\
.15 & .4118 & .276 & 3.622 & 1.48 & 13.25 \\
.10 & .4444 & .372 & 2.690 & 1.813 & 9.937 \\
.05 & .4737 & .515 & 1.941 & 3.15 & 6.46 \\
0 & .5 & 1 & 4 & \\
\bottomrule
\end{tabular}
\end{center}

where the asterisks show the critical values of $\lambda : \Lambda$.

\newpage
\setcounter{page}{5}

It is worth while to transform the equation to new coordinates $X'$, $Y'$, $Z'$ such that $X' = 0$, $Y' = 0$, $Z' = 0$ represent the sides of the triangle formed by the three nodes.

Writing for shortness $X + Y + Z = P$, $YZ + ZX + XY = Q$, $XYZ = R$, the equation is $(q^2P - Q)^2 = 4(2q + 1)PR$.

The expressions of $X$, $Y$, $Z$ in terms of the new coordinates are of the form $X' + \theta P'$, $Y' + \theta P'$, $Z' + \theta P'$, where $P' = X' + Y' + Z'$; writing also $Q' = Y'Z' + Z'X' + X'Y'$, $R' = X'Y'Z'$, then the values of $P$, $Q$, $R$ are $(1 + 3\theta)P'$, $Q' + (2\theta + 3\theta^2)P'^2$, $R' + \theta Q'P' + (\theta^2 + \theta^3)P'^3$.

And the transformed equation is
\begin{multline*}
\bigl[q^2(1 + 3\theta)P'^2 - Q' - (2\theta + 3\theta^2)P'^2\bigr]^2 \\
- 4(2q + 1)(1 + 3\theta)P'\bigl[R' + \theta Q'P' + (\theta^2 + \theta^3)P'^3\bigr] = 0,
\end{multline*}
which is satisfied by $Q' = 0$, $R' = 0$, if only
\[
\bigl[q^2(1 + 3\theta) - 2\theta - 3\theta^2\bigr]^2 = 4(2q + 1)(1 + 3\theta)(\theta^2 + \theta^3),
\]
or, if for a moment $q(1 + 3\theta) = \mu$, the equation is
\[
(\mu^2 - 2\theta - 3\theta^2)^2 = 4(\theta^2 + \theta^3)(2\mu + 1 + 3\theta),
\]
that is,
\[
\mu^4 + \mu^2(-6\theta^2 - 4\theta) + \mu(-8\theta^3 - 8\theta^2) - 3\theta^4 - 4\theta^3 = 0,
\]
that is,
\[
(\mu + \theta)^2(\mu^2 - 2\theta\mu - 3\theta^2 - 4\theta) = 0.
\]

If the new axes pass through the nodes, then $3\theta + 1 = 0$; that is, $q^2(1 + 3\theta) + \theta = 0$, which equation gives the value of $\theta$ for which the new axes have the position in question; substituting in the first instance for $q$ the value $\dfrac{q}{3q+1}$, the equation becomes
\[
[2\theta(1 + \theta)P'^2 + Q']^2 = 4(1 + \theta)P'[\theta^2(1 + \theta)P'^2 + \theta Q'P' + R'],
\]
that is, $Q'^2 = 4(1 + \theta)P'R'$; or, finally, substituting for $\theta$ its value in terms of $q$, the required equation is
\[
(Y'Z' + Z'X' + X'Y')^2 = 4\frac{q+1}{3q+1}X'Y'Z'(X' + Y' + Z').
\]

In particular, for $q = 0$ the equation is
\[
(Y'Z' + Z'X' + X'Y')^2 - 4X'Y'Z'(X' + Y' + Z') = 0,
\]
which is right, since, in the case in question (the tricuspidal curve), we have $X$, $Y$, $Z = X'$, $Y'$, $Z'$.

I remark, in passing, that, taking the equation to be
\[
(Y'Z' + Z'X' + X'Y')^2 = mX'Y'Z'(X' + Y' + Z'),
\]
we may write herein $Z' = 1$, and then the equation takes the form $(Y' + X' + X'Y')^2 = mX'Y'(1 + X' + Y')$; or, what is the same thing, $\sqrt{mX'} + \sqrt{mY'} + \sqrt{X'Y'} + 1 = 0$.

\newpage
\setcounter{page}{6}

The curve may be represented as a unicursal quartic; in fact, taking $\alpha$, $\beta$, $\gamma$ as the roots of
\[
(\alpha - \theta)(\beta - \theta)(\gamma - \theta) = \theta^3 - u\theta^2 + v\theta - w = 0,
\]
we may write
\begin{align*}
x &= \frac{\sqrt{m}\sqrt{(m - 3)^3}}{2\sqrt{m(m - 3)}}\cos \theta + \frac{2(m - 3)}{m}\cos 2\theta,\\
y &= \frac{\sqrt{m}\sqrt{(m - 3)^3}}{2\sqrt{m(m - 3)}}\sin \theta + \frac{2(m - 3)}{m}\sin 2\theta,
\end{align*}
which are the formulae for the description of the trinodal quartic as a unicursal curve.

\newpage
\setcounter{page}{7}

I consider now the second position; viz.\ the horizontal plane now passes through the centre of the tetrahedron, and is parallel to two opposite edges. The equations of the nodal lines are here $(y = 0,\ z = 0)$, $(x = 0,\ z = 0)$, $(x = 0,\ y = 0)$; and if for convenience we assume the distance of the mid-points of opposite edges to be $= 2$, or the half of this $= 1$, then the equations of the faces are
\begin{align*}
X &= x + y + z - 1 = 0,\\
Y &= -x - y + z - 1 = 0,\\
Z &= x - y - z - 1 = 0,\\
W &= -x + y - z - 1 = 0,
\end{align*}
and the equation of the surface is $\sqrt{X} + \sqrt{Y} + \sqrt{Z} + \sqrt{W} = 0$.

Proceeding to rationalise, this is $X + Y + 2\sqrt{XY} = Z + W + 2\sqrt{ZW}$, viz.
\[
2z + 2\sqrt{XY} = 2\sqrt{ZW};
\]
we thence have $4z^2 + 4z\sqrt{XY} + XY = ZW$; or, since $ZW - XY = 4z + 4xy$, this is $4z^2 + 4z\sqrt{XY} + XY = 4z + 4xy + XY$, whence
\[
z + \sqrt{XY} - z = \sqrt{ZW};
\]
\[
(z + xy - z^2)^2 = z^2 - \tfrac{1}{4}(z - 1)^2 - (x^2 - y^2)^2;
\]
\[
2xyz + \tfrac{1}{2}z^2 + \tfrac{1}{2}z(x^2 + y^2) = 0,
\]
a form which puts in evidence the nodal lines.

Considering $z$ as constant, we have the equation of the section; this is a quartic having the node $(x = 0,\ y = 0)$, and two other nodes at infinity on the two axes respectively; moreover, the curve has for bitangents the intersections of its plane with the faces of the tetrahedron; or what is the same thing, attributing to $z$ its constant value, the equations of the bitangents are
\[
x + y + z - 1 = 0,\quad -x - y + z - 1 = 0,\quad x - y - z - 1 = 0,\quad -x + y - z - 1 = 0.
\]

\newpage
\setcounter{page}{8}

Observe that when $c = 1$, we have so that the only real point is $x = 0$, $y = 0$; viz.\ this is a tacnode, having the real tangent $x + y = 0$.

For $c = 0$ (central section) the equation becomes $x^2y^2 = 0$; viz.\ the curve is here the two nodal lines each twice.

It is now easy to trace the changes of form.

\smallskip

$c = 1$; curve is a tacnode, as just mentioned, tangent the dexter diagonal.

\smallskip

$c < 1 > 0$; the tacnode breaks up into two conjugate points on the dexter diagonal, and the curve consists of a loop (enclosed oval) at the origin, and two branches each extending to infinity.

\smallskip

$c = 0$; the loop is reduced to zero, and the two branches are each the axis of $x$ or $y$ taken twice; the only real points of the curve are on the axes.

\smallskip

$c < 0 > -1$; the curve consists of four infinite branches, two on each side of each axis.

\smallskip

$c = -1$; each of the four branches has a cusp at infinity, viz.\ the line infinity is an osculating tangent.

\smallskip

$c < -1 > -3$; the infinite branches have each two (real or imaginary) inflexions; there is a crunode at infinity on each axis.

\smallskip

$c = -3$; the crunodes unite in a triple point at infinity.

\smallskip

$c < -3$; the crunodes disappear (are acnodes), and each of the four branches has again two (real or imaginary) inflexions.

\newpage
\setcounter{page}{9}

The nodal lines of the surface are given by the equations $(y = 0,\ z = 0)$, $(x = 0,\ z = 0)$, $(x = 0,\ y = 0)$. To find the tangent cone at a point on a nodal line, we may proceed as follows.

The general equation of the tangent cone at any point $(x_0, y_0, z_0)$ of a surface $F(x, y, z) = 0$ is obtained by writing $x_0 + \tfrac{X}{W}$, $y_0 + \tfrac{Y}{W}$, $z_0 + \tfrac{Z}{W}$ for $x$, $y$, $z$, and then equating to zero the terms of the lowest order in $X$, $Y$, $Z$, $W$.

For the surface in question, the equation is
\[
F = (z + xy - z^2)^2 - z^2 + \tfrac{1}{4}(z - 1)^2(x^2 - y^2)^2 = 0.
\]

The tangent cone at any point of the nodal line $y = 0$, $z = 0$ (so that $x_0$ is arbitrary) is obtained by writing $x_0 + X$, $Y$, $Z$ for $x$, $y$, $z$; the equation becomes
\[
(Z + x_0Y + XY - Z^2)^2 - Z^2 + \tfrac{1}{4}(Z + x_0^2 - 1)^2(X^2 - Y^2)^2 = 0,
\]
and taking the terms of the lowest order, we have
\[
(Z + x_0Y)^2 - Z^2 + \tfrac{1}{4}(x_0^2 - 1)^2(X^2 - Y^2)^2 = 0.
\]

Equating to zero the discriminant in regard to $Z$, we have
\[
4x_0^2(X^2 - Y^2)^2[(x_0^2 + 1)Y^2 + X^2] - (x_0^2 - 1)^2(X^2 - Y^2)^4 = 0.
\]

That is,
\[
4x_0^2[(x_0^2 + 1)Y^2 + X^2] - (x_0^2 - 1)^2(X^2 - Y^2)^2 = 0;
\]
and substituting herein the values $X = \dfrac{-x_0}{x_0^2 - 1}$ and $Y = \dfrac{1}{x_0^2 - 1}$, we have the equation of the cone, viz.\ this is
\[
\gamma(x^2 + 2y^2 - 1) + 2\gamma(xy + y^2) - \tfrac{1}{4} = 0;
\]
or, what is the same thing,
\[
\gamma(x^2 + 2y^2 - 1) + 2\gamma(xy + y^2) - \tfrac{1}{4} = 0.
\]

\newpage
\setcounter{page}{10}

This is a quadric cone having for its principal planes $y = 0$ (plane through the nodal line and the centre) and $x + y = 0$, $x - y = 0$ (planes through the nodal line and the vertices of the tetrahedron).

The section by the plane $x = 0$ is the conic $2\gamma y^2 + 2\gamma y^2 - \tfrac{1}{4} = 0$, that is $4\gamma y^2 = \tfrac{1}{4}$, or $y = \pm \tfrac{1}{4\sqrt{\gamma}}$; and similarly for the other principal sections. The cone touches the plane $z = 0$ along the line $x + y = 0$.

\newpage
\setcounter{page}{11}

The analytical theory is thus completed for the two positions of the tetrahedron; and we have a complete discussion of the form of the sections by planes parallel to the horizontal plane. The results are in accordance with the model and drawings which I have constructed.

\newpage
\setcounter{page}{12}

The general equation of Steiner's surface may be written in the form
\[
\alpha^2 + \beta^2 + \gamma^2 + \delta^2 - 2\beta\gamma - 2\gamma\alpha - 2\alpha\beta - 2\alpha\delta - 2\beta\delta - 2\gamma\delta = 0,
\]
where $\alpha$, $\beta$, $\gamma$, $\delta$ are linear functions of the coordinates. The four singular tangent planes are $\alpha = 0$, $\beta = 0$, $\gamma = 0$, $\delta = 0$; the three nodal lines are given by the equations $\alpha = \beta = \gamma + \delta$, $\beta = \gamma = \alpha + \delta$, $\gamma = \alpha = \beta + \delta$; and the surface has the property that any plane through a nodal line meets the surface in two conics.

% =====================================================================
% PAPER 557: ON CERTAIN CONSTRUCTIONS FOR BICIRCULAR QUARTICS
% (Pages 33-35 of PDF = pages 13-15)
% =====================================================================

\newpage
\setcounter{page}{13}

\begin{center}
[557\\
]
\end{center}

\section*{\textbf{557.}}
\addcontentsline{toc}{section}{557. On Certain Constructions for Bicircular Quartics}

\begin{center}
{\large \textbf{ON CERTAIN CONSTRUCTIONS FOR BICIRCULAR QUARTICS.}}
\end{center}

\begin{flushright}
{\small [From the Proceedings of the London Mathematical Society, vol.\ v.\ (1873--1874), pp.\ 29--31.\ Read March 12, 1874.]}
\end{flushright}

I call to mind that if $F$, $G$ are any two points and $F'$, $G'$ their antipoints; then the circle on the diameter $FG$ and that on the diameter $F'G'$ are concentric orthotomics, viz.\ they have the same centre, and the sum of the squared radii is $= 0$.

Moreover, if the circles $B$, $B'$ are concentric orthotomics, and the circle $A$ is orthotomic to $B$, then it is a bisector of $B'$, viz.\ it cuts $B'$ at the extremities of a diameter of $B'$; and $B$ is then said to be a bifid circle in regard to $A$.

Given two real circles, these have an axial orthotomic, the circle, centre on the line of centres at its intersection with the radical axis, which cuts at right angles the given circles; viz.\ this axial orthotomic is real if the circles have no real intersection; but if the intersections are real, then the axial orthotomic is a pure imaginary.

The above remarks have an obvious application to the theory of bicircular quartics; viz.\ a bicircular quartic is the envelope of a variable circle, having its centre on a conic, and orthotomic to a circle: it may be that this circle is a pure imaginary. We then replace it by the concentric orthotomic, and say that the curve is the envelope of a variable circle having its centre on a conic and bisecting a circle. We have thus a real form for cases which originally present themselves under an imaginary form.

\newpage
\setcounter{page}{14}

\textbf{The Bicircular Quartic with given vertices.}

First, if the vertices are all real, say $f$, $g$, $h$, $k$ (these being the foci or double foci of the curve); then taking $F$, $G$, $H$, $K$ to be the points whose coordinates are $f$, $g$, $h$, $k$ respectively, and $F'$, $G'$ the antipoints of $F$, $G$, and $H'$, $K'$ the antipoints of $H$, $K$; we have the two concentric orthotomic circles $F'G'$ and $H'K'$.

The axial orthotomic of these two circles is real (since the circles have no real intersection), and calling this the circle $\Omega$, the curve is the envelope of a variable circle having its centre on a conic and orthotomic to $\Omega$.

To obtain the conic, we remark that the vertices $f$, $g$ belong to the circle $F'G'$, and $h$, $k$ to the circle $H'K'$; and if on the line of centres (which is also the radical axis of the two circles) we describe a conic having $F'G'$ and $H'K'$ as its director circles, then the variable circle is orthotomic to $\Omega$, and its centre describes the conic. The conic is in fact a central conic, the axes of which are the line of centres and the perpendicular thereto through the centre.

If the conic is an ellipse, the curve is a simple oval; if it is a hyperbola, the curve consists of two ovals.

\newpage
\setcounter{page}{15}

Secondly, if the vertices are two real, two imaginary, say $f$, $g = \alpha \pm \beta i$; $h$, $k$, we modify the first or third construction; viz.\ if $F'$, $G'$ are the antipoints of $F$, $G$; then on $F'G'$ as diameter describe a circle, and on $HK$ as diameter a circle. On the line terminated by the two centres (as transverse or conjugate axis) describe a conic $\Theta_1$, and describe the axial bisector-orthotomic circle $\Omega_1$ of the two circles; viz.\ this is the circle (centre on the axis of symmetry) which bisects the circle $F'G'$, and cuts at right angles the circle $HK$; then the curve is the envelope of the variable circle having its centre on $\Theta_1$ and orthotomic to $\Omega_1$.

Thirdly, if the vertices are all imaginary, say $f$, $g = \alpha \pm \beta i$; $h$, $k = \gamma \pm \delta i$, we modify the first or third construction in like manner.

The curve is in this case quadrantal: having, besides the original axis of symmetry, another axis of symmetry through $O$, at right angles thereto.

% =====================================================================
% PAPER 558: A GEOMETRICAL INTERPRETATION (Pages 36-37 of PDF = pages 16-17)
% =====================================================================

\newpage
\setcounter{page}{16}

\begin{center}
[558\\
]
\end{center}

\section*{\textbf{558.}}
\addcontentsline{toc}{section}{558. A Geometrical Interpretation of the Equations Obtained by Equating to Zero the Resultant and the Discriminants of Two Binary Quantics}

\begin{center}
{\large \textbf{A GEOMETRICAL INTERPRETATION OF THE EQUATIONS}}\\[4pt]
{\large \textbf{OBTAINED BY EQUATING TO ZERO THE RESULTANT AND THE}}\\[4pt]
{\large \textbf{DISCRIMINANTS OF TWO BINARY QUANTICS.}}
\end{center}

\begin{flushright}
{\small [From the Proceedings of the London Mathematical Society, vol.\ v.\ (1873--1874), pp.\ 31--33.\ Read March 12, 1874.]}
\end{flushright}

Consider the equations $U = (a, b, \ldots \wr x, 1)^\lambda = 0$, $U' = (a', b', \ldots \wr x, 1)^{\lambda'} = 0$; and equating to zero the discriminants of the two functions respectively, and also the resultant of the two functions, let the equations thus obtained be
\[
\Delta = (a, b, \ldots)^{\lambda-2} = 0,\quad \Delta' = (a', b', \ldots)^{\lambda'-2} = 0,\quad R = (a, b, \ldots)^\lambda (a', b', \ldots)^{\lambda'} = 0.
\]

I take $(a, b, \ldots)$, $(a', b', \ldots)$ to be linear functions of the coordinates $(x, y, z)$; and $t$ to be an indeterminate parameter. Hence $U = 0$ represents a line the envelope whereof is the curve $\Delta = 0$, or, what is the same thing, the equation $U = 0$ represents any tangent of the curve $\Delta = 0$; this is a unicursal curve of the order $2\lambda - 2$ and class $\lambda$, with $3(\lambda - 2)$ cusps and $\tfrac{1}{2}(\lambda - 2)(\lambda - 3)$ nodes.

Similarly $U' = 0$ represents a line the envelope whereof is the curve $\Delta' = 0$, a unicursal curve of the order $2\lambda' - 2$ and class $\lambda'$, with $3(\lambda' - 2)$ cusps and $\tfrac{1}{2}(\lambda' - 2)(\lambda' - 3)$ nodes.

\newpage
\setcounter{page}{17}

The envelope of the line which is the join of the corresponding points of contact of the two curves is the curve $R = 0$, a unicursal curve of the order $\lambda + \lambda'$, with $\tfrac{1}{2}(\lambda + \lambda' - 1)(\lambda + \lambda' - 2)$ nodes and no cusps; consequently of the class $2(\lambda + \lambda' - 1)$.

It is to be shown that the curve $R = 0$ touches the curve $\Delta = 0$ in $\lambda' + 2\lambda - 2$ points, and similarly the curve $\Delta' = 0$ in $2\lambda' + \lambda - 2$ points.

In fact, consider any tangent $T'$ of the curve $\Delta'$; let this meet the curve $\Delta$ in a point $A$, and let $Q$ be the tangent at $A$ to the curve $\Delta$; suppose, moreover, that $T$ is the tangent of $\Delta$ corresponding to the tangent $T'$ of $\Delta'$. Then if $Q$ and $T$ coincide, the corresponding tangent of $T'$ will be $Q$, and the curve $R$ will pass through $A$. It is easy to see that in this case the curves $R$, $\Delta$ will touch at $A$.

Again, if $P$ be a tangent from $A$ to the curve $\Delta$, then, if $P$ and $T$ coincide, the corresponding tangent of $T'$ will be $P$, and the curve $R$ will pass through $A$; but in this case the point $A$ will be a mere intersection, not a point of contact, of the two curves.

Hence the points of contact of $R$ with $\Delta$ correspond to the tangents $T'$ which touch $\Delta$, that is, to the $2(\lambda - 1)$ tangents from the cusps of $\Delta'$ to the curve $\Delta$, together with the $\lambda'$ common tangents of $\Delta$ and $\Delta'$; in all $2(\lambda - 1) + \lambda' = \lambda' + 2\lambda - 2$ points of contact. And similarly the number of points of contact of $R$ with $\Delta'$ is $2\lambda' + \lambda - 2$.

% =====================================================================
% PAPER 559: [NOTE ON INVERSION] (Page 38 of PDF = page 18)
% =====================================================================

\newpage
\setcounter{page}{18}

\begin{center}
[559\\
]
\end{center}

\section*{\textbf{559.}}
\addcontentsline{toc}{section}{559. [Note on Inversion]}

\begin{center}
{\large \textbf{[NOTE ON INVERSION.]}}
\end{center}

\begin{flushright}
{\small [From the Proceedings of the London Mathematical Society, vol.\ v.\ (1873--1874), p.\ 112.]}
\end{flushright}

The inverse of the anchor ring (in the foregoing paper called the cyclide) is in fact the general binodal cyclide or binodal bicircular quartic; viz.\ assuming it to be a cyclide (bicircular quartic), to see that it is binodal, it need only be observed that the anchor ring is binodal (has two real or imaginary conic points, viz.\ these are the intersections of the circles in the several axial planes); and to see that it is the general binodal cyclide, we have only to count the constants; viz.\ the general cyclide or surface
\begin{multline*}
(x^2 + y^2 + z^2)^2 + (x^2 + y^2 + z^2)(\alpha x + \beta y + \gamma z)\\
+ (a, b, c, d, f, g, h, l, m, n)(x, y, z, 1)^2 = 0
\end{multline*}
contains 13 constants, and therefore the binodal cyclide $13 - 2 = 11$ constants. But the anchor ring, irrespective of position, contains 2 constants; centre of inversion, taken in given axial plane, has 2 constants; radius of inversion, 1 constant; in all $2 + 2 + 1 = 5$ constants for the form, and 6 for the position, in all 11 constants, the same number as for the general binodal cyclide.

% =====================================================================
% PAPER 560: ADDITION TO LORD RAYLEIGH'S PAPER (Pages 39-40 of PDF = pages 19-20)
% =====================================================================

\newpage
\setcounter{page}{19}

\begin{center}
[560\\
]
\end{center}

\section*{\textbf{560.}}
\addcontentsline{toc}{section}{560. [Addition to Lord Rayleigh's Paper ``On the Numerical Calculation of the Roots of Fluctuating Functions'']}

\begin{center}
{\large \textbf{[ADDITION TO LORD RAYLEIGH'S PAPER ``ON THE NUMERICAL}}\\[4pt]
{\large \textbf{CALCULATION OF THE ROOTS OF FLUCTUATING FUNCTIONS.'']}}
\end{center}

\begin{flushright}
{\small [From the Proceedings of the London Mathematical Society, vol.\ v.\ (1873--1874), pp.\ 123, 124.\ November 22.]}
\end{flushright}

Prof.\ Cayley, to whom Lord Rayleigh's paper was referred, pointed out that a similar result may be attained by a method given in a paper by Encke, ``Allgemeine Aufloesung der numerischen Gleichungen,'' Crelle, t.\ xxii.\ (1841), pp.\ 193--248, as follows:

Taking the equation
\[
0 = 1 - ax + bx^2 - cx^3 + dx^4 - ex^5 + fx^6 - gx^7 + hx^8 - \ldots;
\]
if the equation whose roots are the squares of these is
\[
0 = 1 - a_1 x + b_1 x^2 - c_1 x^3 + \ldots,
\]
then
\begin{align*}
a_1 &= a^2 - 2b,\\
b_1 &= b^2 - 2ac + 2d,\\
c_1 &= c^2 - 2bd + 2ae - 2f,\\
d_1 &= d^2 - 2ce + 2bf - 2ag + 2h,
\end{align*}
\&c.; and we may in the same way derive $a_2$, $b_2$, $c_2$, \&c.\ from $a_1$, $b_1$, $c_1$, \&c., and so on.

As regards the function
\[
\phi^n(x) = \frac{2^n \cdot n!}{x^n}\left\{J_n(x)\right\} = 1 - \frac{x^2}{2 \cdot 2n + 2} + \frac{x^4}{2 \cdot 4 \cdot 2n + 2 \cdot 2n + 4} - \ldots,
\]
we have as follows:

\newpage
\setcounter{page}{20}

\begin{align*}
a^{-1} &= 2^1 \cdot n + 1,\\
b^{-1} &= 2^2 \cdot n + 1 \cdot n + 2,\\
c^{-1} &= 2^3 \cdot 3 \cdot n + 1 \ldots n + 3,\\
d^{-1} &= 2^4 \cdot 3 \cdot n + 1 \ldots n + 4,\\
e^{-1} &= 2^5 \cdot 3 \cdot 5 \cdot n + 1 \ldots n + 5,\\
f^{-1} &= 2^6 \cdot 3^2 \cdot 5 \cdot n + 1 \ldots n + 6,\\
g^{-1} &= 2^7 \cdot 3^2 \cdot 5 \cdot 7 \cdot n + 1 \ldots n + 7,\\
h^{-1} &= 2^8 \cdot 3^2 \cdot 5 \cdot 7 \cdot n + 1 \ldots n + 8,\\
\intertext{and}
a_1^{-1} &= 2^4 \cdot (n + 1)^2 \cdot n + 2,\\
b_1^{-1} &= 2^7 \cdot (n + 1 \cdot n + 2)^2 \cdot n + 3 \cdot n + 4,\\
c_1^{-1} &= 2^{10} \cdot 3 \cdot (n + 1 \ldots n + 3)^2 \cdot n + 4 \ldots n + 6,\\
d_1^{-1} &= 2^{13} \cdot 3 \cdot (n + 1 \ldots n + 4)^2 \cdot n + 5 \ldots n + 8,\\
a_2 &= \frac{5n^2 + 23n + 42}{2^8 \cdot (n + 1)^2 \cdot (n + 2)^2 \cdot n + 3 \cdot n + 4},\\
b_2 &= \frac{429n^4 + 7640n^3 + 53702n^2 + 185430n + 311387n + 202738}{2^{16} \cdot (n + 1)^8 \cdot (n + 2)^4 \cdot (n + 3 \cdot n + 4)^2 \cdot n + 5 \cdot n + 6 \cdot n + 7 \cdot n + 8}.
\end{align*}

If $n = 0$, whence
\[
p_1 = a_2 = \frac{101369}{2^{15} \cdot 3^2 \cdot 5 \cdot 7} = p_1 \text{ suppose;}
\]
\[
p_1 = 2.404825.
\]

[The quantities $p_1$, $p_2$, \ldots\ are the roots of the function $J_n(x)$ in increasing order of magnitude, so that, as these roots are all real, it follows that for $J_0(x)$,
\[
a = 2p_1^{-2},\quad a_1 = 2p_1^{-4},\quad a_2 = 2p_1^{-6},\quad a_3 = 2p_1^{-8},\quad \ldots]
\]

% =====================================================================
% PAPER 561: ON THE GEOMETRICAL REPRESENTATION OF CAUCHY'S THEOREMS
% OF ROOT-LIMITATION (Pages 41-50 of PDF = pages 21-30)
% =====================================================================

\newpage
\setcounter{page}{21}

\begin{center}
[561\\
]
\end{center}

\section*{\textbf{561.}}
\addcontentsline{toc}{section}{561. On the Geometrical Representation of Cauchy's Theorems of Root-Limitation}

\begin{center}
{\large \textbf{ON THE GEOMETRICAL REPRESENTATION OF CAUCHY'S}}\\[4pt]
{\large \textbf{THEOREMS OF ROOT-LIMITATION.}}
\end{center}

\begin{flushright}
{\small [From the Transactions of the Cambridge Philosophical Society, vol.\ xii.\ Part ii.\ (1877), pp.\ 395--413.\ Read February 16, 1874.]}
\end{flushright}

There is contained in Cauchy's Memoir ``Calcul des Indices des Fonctions,'' Journ.\ de l'Ecole Polytech.\ t.\ xv.\ (1837) a general theorem, which, though including a well-known theorem in regard to the imaginary roots of a numerical equation, seems itself to have been almost lost sight of.

In the general theorem (say Cauchy's two-curve theorem) we have in a plane two curves $P = 0$, $Q = 0$, and the real intersections of these two curves, or say the ``roots,'' are divided into two sets according as the Jacobian $\partial_x P \cdot \partial_y Q - \partial_x Q \cdot \partial_y P$ is positive or negative, say these are the Jacobian-positive and the Jacobian-negative roots: and the question is to determine for the roots within a given contour or circuit, the difference of the numbers of the roots belonging to the two sets respectively.

I reproduce the whole theory, presenting it under a completely geometrical form, viz.\ I establish between the two sets of roots the distinction of right- and left-handed: and (availing myself of a notion due to Prof.\ Sylvester) I give a geometrical form to the theoretic rule, making it depend on the ``intercalation'' of the intersections of the two curves with the circuit: I also complete the Sturmian process in regard to the sides of the rectangle: the memoir contains further researches in regard to the theory.

\textbf{1. Right- and Left-handed Roots.}

Consider any point of intersection of the two curves. There will be about this point four regions, sable and argent being opposite to each other, as also gules and azure; whence selecting an order sable, gules, argent, azure; if to have the colours in this order we have to go about the point, or root, right-handedly, the root is right-handed: but if to have the colours in this order we have to go about the root left-handedly, the root is left-handed.

It is easy to see that if the Jacobian is positive, the root is right-handed; and if the Jacobian is negative, the root is left-handed.

\textbf{2. The Circuit.}

I consider in general a circuit not containing within it any root; as a simple example let the circuit lie altogether in a sable region: the intercalation is here $= 0$, and the circuit is said to be neutral.

If the circuit contains within it a single root, and if this is right-handed, the intercalation is $+P - Q$; and the Index of the circuit, = number of right-handed roots minus number of left-handed roots, is $= +1$.

\newpage
\setcounter{page}{22}

\textbf{3. Intercalation-theory for a right line.}

Considering then the case where the trajectory is a line parallel to the axis of $X$, $P$ and $Q$ will denote given rational functions of $x$; the curves $P = 0$, $Q = 0$ being of course each of them a set of right lines parallel to the axis of $y$: the regions will be bands each of them included between two consecutive lines $P = 0$, or between two consecutive lines $Q = 0$.

A band included between two consecutive lines $P = 0$ is a $P$-band; and so a band included between two consecutive lines $Q = 0$ is a $Q$-band. In regard to a $P$-band, the signs of $P$ on the two sides thereof are + and -, or - and +; we may distinguish these as the +$P$ side and the -$P$ side respectively: and the like as regards a $Q$-band.

\textbf{4. Bicolumns.}

We have a bicolumn: that is, a double column of signs; say the first or $P$-column, consisting of the signs of $P$ at the several roots of $Q = 0$; and the second or $Q$-column, consisting of the signs of $Q$ at the several roots of $P = 0$; the roots are here taken in order of increasing magnitude, and the bicolumn is read downwards.

The intercalation of the bicolumn is derived from the two columns taken separately; the $P$-column contributes only $P$'s, the $Q$-column contributes only $Q$'s: the two letters are written down in the order in which they present themselves; the signs + and - are attributed as above: and we have thus the intercalation.

\newpage
\setcounter{page}{23}

\textbf{5. Gain of a bicolumn.}

The gain of a bicolumn is the number of the letters of the intercalation taken with the proper signs. The rule for the signs is: for the first letter the sign is +, and the signs are alternately + and -: but the signs of the $P$'s and the signs of the $Q$'s are each of them independently + and - according to the sign of the letter in the intercalation.

As regards the final result, a bicolumn may be replaced by its gain, or single letter with its sign.

\textbf{6. Decomposition of a bicolumn.}

A bicolumn may be divided in any manner into parts, taking always the last row of any part as being also the first row of the next succeeding part. This being so, the gain of the whole bicolumn is equal to the sum of the gains of its parts.

In a bicolumn of two rows, if we reverse either row (that is, write therein $-$ for $+$ and $+$ for $-$), we reverse the signs of the letters contributed by the reversed column, and leave the gain unaltered---of course the intercalation (considered irrespectively of sign) is in each case unaltered.

It is convenient to take the signs in such manner that for the initial value of $x$, the signs of $P$, $Q$ shall be each positive: or, what is the same thing, taking $P$, $Q$ with their proper signs, we may in the bicolumn, by reversing if necessary, take the signs of the first row each positive.

\textbf{7. Application to Sturm's theorem.}

For a rectangle the sides of which are parallel to the axes, the index is equal to the gain of the bicolumn formed by the signs at the corners: viz.\ reading the corners in order, beginning with the left-hand corner of the under side, and going round right-handedly, we have a bicolumn of four rows: and the index is equal to the gain of this bicolumn.

\newpage
\setcounter{page}{24}

The index is thus equal to the sum of the gains of the two side-bicolumns and the two terminal columns; and we have thence the expression for the index in the form required for the application of Sturm's theorem.

\textbf{8.}

Consider a circuit not containing within it any root; as a simple example let the circuit lie altogether in a sable region: the intercalation is here $= 0$, and the circuit is said to be neutral.

If the circuit contains within it a single root, and if this is right-handed, the intercalation is $+P - Q$; and the Index of the circuit, = number of right-handed roots minus number of left-handed roots, is $= +1$.

If the circuit contains within it any number of right-handed roots, and any number of left-handed roots, then the Index is equal to the number of right-handed roots less the number of left-handed roots.

\textbf{9.}

In the case of a circuit not containing any root, the intercalation is either $(+P - Q)$, say this is a positive circuit, or else $(-P + Q)$, say this is a negative circuit. There is of course the neutral circuit $(PQ)$ for which the intercalation vanishes.

\textbf{10.}

Consider a circuit not containing within it any root; as a simple example let the circuit lie altogether in a sable region: the intercalation is here $= 0$, and the circuit is said to be neutral.

If the circuit contains within it a single root, and if this is right-handed, the intercalation is $+P - Q$; and the Index of the circuit, = number of right-handed roots minus number of left-handed roots, is $= +1$.

\newpage
\setcounter{page}{25}

\textbf{11.}

Consider a circuit not containing within it any root; as a simple example let the circuit lie altogether in a sable region: the intercalation is here $= 0$, and the circuit is said to be neutral.

If the circuit contains within it a single root, and if this is right-handed, the intercalation is $+P - Q$; and the Index of the circuit, = number of right-handed roots minus number of left-handed roots, is $= +1$.

\textbf{12.}

In the second case, that is, when the intercalations are contrary, they counteract each other in forming the intercalation of the circuit: it is the difference of the numbers of letters, or twice the difference of the indices, which is evenly even, and the half of this, or the difference of the indices, which is the index of the circuit: one intercalation contains a number of letters greater by an odd number than the number of letters in the other intercalation, and the index of the circuit is the half of an odd number, that is, it is a fraction with denominator 2, or say it is a fractional index: this can only happen when the circuit passes through a root.

\textbf{13.}

Suppose $CA$ similar to $ABC$, and $AG$ similar to $CDA$; then $ABCA$ is also similar to $ABC$, and $AGDA$ similar to $CDA$; viz.\ $ABC$, $GA$ and $ABCA$ are similar to each other, and contrary to $AC$, $CDA$, $ACDA$ which are also similar to each other.

Also
\begin{align*}
\text{Ind.}\ ABCDA &= \text{Ind.}\ ABC + \text{Ind.}\ CDA,\\
\text{Ind.}\ ABCA &= \text{Ind.}\ ABC + \text{Ind.}\ CA,\\
\text{Ind.}\ ACDA &= \text{Ind.}\ CDA + \text{Ind.}\ AC,
\end{align*}
and thence
\[
\text{Ind.}\ ABCDA = \text{Ind.}\ ABCA + \text{Ind.}\ ACDA.
\]

\newpage
\setcounter{page}{26}

\textbf{14.}

The theory of intercalations as thus applied to a polygon gives a complete geometrical representation of Cauchy's theorem for any contour whatever. For a rectangle the sides of which are parallel to the axes, the index is equal to the gain of the bicolumn formed by the signs at the corners: viz.\ reading the corners in order, beginning with the left-hand corner of the under side, and going round right-handedly, we have a bicolumn of four rows: and the index is equal to the gain of this bicolumn.

The index is thus equal to the sum of the gains of the two side-bicolumns and the two terminal columns; and we have thence the expression for the index in the form required for the application of Sturm's theorem.

\textbf{15.}

Consider a circuit not containing within it any root; as a simple example let the circuit lie altogether in a sable region: the intercalation is here $= 0$, and the circuit is said to be neutral.

If the circuit contains within it a single root, and if this is right-handed, the intercalation is $+P - Q$; and the Index of the circuit, = number of right-handed roots minus number of left-handed roots, is $= +1$.

\newpage
\setcounter{page}{27}

\textbf{16.}

Suppose $CA$ similar to $ABC$, and $AG$ similar to $CDA$; then $ABCA$ is also similar to $ABC$, and $AGDA$ similar to $CDA$; viz.\ $ABC$, $GA$ and $ABCA$ are similar to each other, and contrary to $AC$, $CDA$, $ACDA$ which are also similar to each other.

Also
\begin{align*}
\text{Ind.}\ ABCDA &= \text{Ind.}\ ABC + \text{Ind.}\ CDA,\\
\text{Ind.}\ ABCA &= \text{Ind.}\ ABC + \text{Ind.}\ CA,\\
\text{Ind.}\ ACDA &= \text{Ind.}\ CDA + \text{Ind.}\ AC,
\end{align*}
and thence
\[
\text{Ind.}\ ABCDA = \text{Ind.}\ ABCA + \text{Ind.}\ ACDA.
\]

\textbf{17.}

In the case where the circuit is a polygon, and most easily when it is a rectangle the sides of which are parallel to the two axes respectively, the excess in question can be actually determined by means of an application of Sturm's theorem successively to each side of the polygon, or rectangle.

\newpage
\setcounter{page}{28}

\textbf{18. Intercalation-theory for a right line.}

Considering then the case where the trajectory is a line parallel to the axis of $X$, $P$ and $Q$ will denote given rational functions of $x$; the curves $P = 0$, $Q = 0$ being of course each of them a set of right lines parallel to the axis of $y$: the regions will be bands each of them included between two consecutive lines $P = 0$, or between two consecutive lines $Q = 0$.

A band included between two consecutive lines $P = 0$ is a $P$-band; and so a band included between two consecutive lines $Q = 0$ is a $Q$-band. In regard to a $P$-band, the signs of $P$ on the two sides thereof are + and -, or - and +; we may distinguish these as the +$P$ side and the -$P$ side respectively: and the like as regards a $Q$-band.

\textbf{19.}

We have a bicolumn: that is, a double column of signs; say the first or $P$-column, consisting of the signs of $P$ at the several roots of $Q = 0$; and the second or $Q$-column, consisting of the signs of $Q$ at the several roots of $P = 0$; the roots are here taken in order of increasing magnitude, and the bicolumn is read downwards.

\textbf{20.}

The intercalation of the bicolumn is derived from the two columns taken separately; the $P$-column contributes only $P$'s, the $Q$-column contributes only $Q$'s: the two letters are written down in the order in which they present themselves; the signs + and - are attributed as above: and we have thus the intercalation.

\newpage
\setcounter{page}{29}

\textbf{21. A bicolumn may be divided in any manner.}

A bicolumn may be divided in any manner into parts, taking always the last row of any part as being also the first row of the next succeeding part. This being so, the gain of the whole bicolumn is equal to the sum of the gains of its parts.

In a bicolumn of two rows, if we reverse either row (that is, write therein $-$ for $+$ and $+$ for $-$), we reverse the signs of the letters contributed by the reversed column, and leave the gain unaltered---of course the intercalation (considered irrespectively of sign) is in each case unaltered.

It is convenient to take the signs in such manner that for the initial value of $x$, the signs of $P$, $Q$ shall be each positive: or, what is the same thing, taking $P$, $Q$ with their proper signs, we may in the bicolumn, by reversing if necessary, take the signs of the first row each positive.

\textbf{22.}

The gain of a bicolumn is the number of the letters of the intercalation taken with the proper signs. The rule for the signs is: for the first letter the sign is +, and the signs are alternately + and -: but the signs of the $P$'s and the signs of the $Q$'s are each of them independently + and - according to the sign of the letter in the intercalation.

As regards the final result, a bicolumn may be replaced by its gain, or single letter with its sign.

\textbf{23.}

We may, in regard to the final result, replace any three consecutive rows by a single row containing the product of the signs: the bicolumn is thus reduced to a bicolumn of two rows, the gain of which is equal to the gain of the original bicolumn.

\newpage
\setcounter{page}{30}

\textbf{24.}

The application to Sturm's theorem is immediate. For a rectangle the sides of which are parallel to the axes, the index is equal to the gain of the bicolumn formed by the signs at the corners: viz.\ reading the corners in order, beginning with the left-hand corner of the under side, and going round right-handedly, we have a bicolumn of four rows: and the index is equal to the gain of this bicolumn.

The index is thus equal to the sum of the gains of the two side-bicolumns and the two terminal columns; and we have thence the expression for the index in the form required for the application of Sturm's theorem.

\vspace{20pt}

\begin{center}
\rule{0.5\textwidth}{0.4pt}
\end{center}

\vspace{12pt}

\noindent\textit{[End of Papers 556--561, covering pages 1--30 of Volume IX, corresponding to pages 1--50 of the PDF source.]}

\vspace{12pt}

\noindent\textit{Transcription note: The source PDF contains significant OCR degradation. Mathematical formulae have been reconstructed from context. The original papers continue beyond this point in the source volume.}

\end{document}
